\answer{OP37}
将点排序为 $X_{(1)}<\cdots<X_{(N)}$，令内部间距
\[
D_i=X_{(i+1)}-X_{(i)},\qquad 1\le i\le N-1.
\]
互为最近邻的两点必定相邻。对每一条相邻边，令 $I_i$ 表示两端互为最近邻。端点边 $i=1$ 只需满足 $D_1<D_2$，末端边同理；内部边 $2\le i\le N-2$ 则需满足
\[
D_i<D_{i-1},\qquad D_i<D_{i+1}.
\]
均匀次序统计量的所有间距（连同两端边界间距）交换对称，且相等概率为零。因此端点边成为配对的概率为 $1/2$，内部边的三个间距中自身最小的概率为 $1/3$。

每个点只有一个最近邻，所以不同互选配对没有公共端点；每个配对恰贡献两个好点。若 $G$ 为好点数，则
\[
G=2\sum_{i=1}^{N-1}I_i,
\quad
\E G=2\left(\frac12+\frac12+\frac{N-3}{3}\right)
=\boxed{\frac{2N}{3}}.
\]
$N=3$ 时内部边求和为空，上式仍成立。这里数的是点；互选配对数的期望为 $N/3$。
